\scanpage{9}
% The large diagonal strokes cancel the argument through the last scalar equation.
\deleted{Plaçons nous sur corps $k$, et supposons $\varepsilon\ne1$. (donc car. $\ne2$) -- le cas $\varepsilon=1$ étant le cas déjà traité (sous \unsure{le|la} \unsure{restriction} que $\beta,\rho$ non nuls tous deux). On doit donc avoir \uline{$\rho$ inv.} par la première équation, donc $a=\frac{\deleted{\unsure{\varepsilon|1-\varepsilon}}2}{\rho}\pi=\lambda\pi$, \deleted{donc} \inserted{$(\lambda=\frac2\rho)$} la $2^{\mathrm{ième}}$ équation donne $\rho\lambda=2\deleted{\unsure{-\varepsilon}}$ \deleted{déjà} satisfaite, enfin la dernière s'écrit}
\[
\begin{aligned}
&\deleted{\lambda^2\underbrace{f(\pi)}_{-\widetilde\rho\widetilde\sigma}\pi'\unsure{-|+}\deleted{\varepsilon}\lambda\underbrace{\psi(\pi)}_{-\sigma\pi'}=0}\\
\deleted{\text{i.e.}}\quad&\deleted{\lambda(\deleted{\unsure{-}}\underbrace{\widetilde\rho\widetilde\sigma\lambda}_{\frac{2\widetilde\rho\widetilde\sigma}{\rho}=\frac{\rho\sigma}{\rho}=\sigma}\unsure{-|+}\deleted{\varepsilon}\sigma)=0}
\end{aligned}
\]
\deleted{ce qui exige, puisque $\lambda\ne0$, que}
\[
\deleted{\widetilde\rho\widetilde\sigma\lambda\unsure{-|+}\deleted{\varepsilon}\sigma=0}
\]
\deleted{ou encore}
\[
\begin{aligned}
&\deleted{\unsure{\rho\,\tfrac{\widetilde\sigma}{2}(1-\varepsilon)}}\qquad\deleted{\sigma=\widetilde\rho\widetilde\sigma}\\
&\deleted{(\unsure{1-\varepsilon})\frac{2\widetilde\rho\widetilde\sigma}{\rho}+\varepsilon\sigma=0}
\end{aligned}
\]

Sur un anneau intègre, les seuls cas pour $\varepsilon$ sont $\varepsilon=+1$ (déjà étudié, sous la restriction que $\rho,\beta$ non tous deux inv) ou $\varepsilon=-1,\ne+\uline{1}$ (donc \deleted{car. $\ne2$)} corps des fractions $K$ de car $\ne2$) donc $\rho a=2\pi$ donc $a=\frac2\rho\pi$, $\deleted{\varepsilon}u=-\mathrm{id}+\mathord{?}\pi'\otimes\frac2\rho\pi$, c'est l'antipodisme, \inserted{\unsure{voir le Signe|vérifier le Signe}} dont on sait bien qu'il conserve $f$, et \deleted{\unsure{conserve|inverse}} $\pi,\pi'$ .. -- Donc

\uline{Proposition} \deleted{\unsure{Prop}} \inserted{(Corps de base)} Si $\beta$ ou $\rho$ est inv. \deleted{et car. $\ne2$,} alors le seul autom de $(E,f,\pi,\pi')$ qui induise l'identité

% archive sheet number: 8

